
\section{Potencjały Bessela i przestrzenie \texorpdfstring{$H^s$}{Hs}}
\label{sec:dt-bessel}
\providecommand{\atlasref}[2]{#2}

Ujednolicamy normalizację przestrzeni Bessela z rozdziałem
o dystrybucjach temperowanych. Por.
\atlasref{PDE-08::thm:pde8-Hs}{RRC-08: potencjały Bessela}.


Dla $s\in\R$ połóżmy $b_s(\xi)=(1+|\xi|^2)^{s/2}$.
Najpierw uzasadnimy, że $b_s$ jest dopuszczalnym
mnożnikiem dystrybucji temperowanych.
Każda jego pochodna jest skończoną sumą wyrazów
postaci
\[
 P(\xi)(1+|\xi|^2)^{s/2-k},
 \qquad k\in\N_0,
\]
gdzie $P$ jest wielomianem. Dla pochodnej rzędu
zero mamy $P=1$, $k=0$. Różniczkowanie takiego
wyrazu względem $\xi_j$ daje
\[
 (\partial_j P)b_{s-2k}
 +(s-2k)\xi_j P b_{s-2k-2},
\]
co uzasadnia indukcję.
Dla dowolnego rzeczywistego $a$ mamy
$(1+|\xi|^2)^{a/2}\le(1+|\xi|)^{\max\{a,0\}}$:
dla $a\le0$ lewa strona nie przekracza $1$,
a dla $a>0$ korzystamy z
$\sqrt{1+|\xi|^2}\le1+|\xi|$.
Każda pochodna $b_s$ ma zatem wzrost co najwyżej
wielomianowy (można zwiększyć wykładnik do
liczby całkowitej). Na mocy twierdzenia~\ref{thm:tempered-operators22}
o mnożnikach dystrybucji temperowanych
mnożenie przez $b_s$ działa na $\mathcal S$
i na $\mathcal S'$.

\begin{definicja}[Przestrzeń Sobolewa typu Bessela]
\label{def:dt-Hs}
Dla dowolnego $s\in\R$ definiujemy
\[
 H^s(\R^d)=\bigl\{u\in\mathcal S'(\R^d):
 b_s\widehat u=T_h\text{ dla pewnego }
 h\in L^2(\R^d)\bigr\}.
\]
Z liniowości transformaty i mnożenia przez $b_s$
oraz liniowości odwzorowania $h\mapsto T_h$
wynika, że $H^s$ jest zespoloną podprzestrzenią
wektorową $\mathcal S'$.
Funkcja $h$ jest określona jednoznacznie jako
klasa $L^2$ dzięki iniektywności odwzorowania
funkcji lokalnie całkowalnych w dystrybucje.
Definiujemy
$\|u\|_{H^s}=\|h\|_2$.
Mnożąc równość $b_s\widehat u=T_h$ przez
$b_{-s}$, otrzymujemy dla $\varphi\in\mathcal S$
\[
 \langle\widehat u,\varphi\rangle
 =\langle T_h,b_{-s}\varphi\rangle
 =\int_{\R^d}b_{-s}(\xi)h(\xi)\varphi(\xi)
 \,\mathrm d\lambda_d(\xi).
\]
Całka jest bezwzględnie zbieżna z nierówności
Cauchy'ego--Schwarza, ponieważ $h\in L^2$
i $b_{-s}\varphi\in\mathcal S\subset L^2$.
Zatem $\widehat u=T_{b_{-s}h}$ i możemy zapisać
\[
 \|u\|_{H^s}^2=\int_{\R^d}
 (1+|\xi|^2)^s|\widehat u(\xi)|^2\,
 \mathrm d\lambda_d(\xi),
\]
gdzie pod całką $\widehat u$ oznacza funkcję
$b_{-s}h$, lokalnie należącą do $L^2$.
Iloczyn jest lokalnie kwadratowalny, ponieważ
$b_{-s}$ jest ograniczony na każdym zbiorze zwartym.
Operator
\[
 J^s u=(I-\Delta)^{s/2}u
 :=\mathcal F^{-1}(b_s\widehat u)
\]
określamy na $\mathcal S'$ przez mnożnik Fouriera,
a nie przez formalne potęgowanie operatora
różniczkowego.
Dla ujemnego $s$ element $H^s$ może być
dystrybucją nieregularną; wymóg $L^2$ dotyczy
ważonej transformaty, a nie samego $u$.
\end{definicja}

Z równości $b_sb_t=b_{s+t}$ wynika
$J^sJ^t=J^{s+t}$ i $J^0=I$, a więc
$(J^s)^{-1}=J^{-s}$ na $\mathcal S'$.
Istotnie, mnożenie dystrybucji przez dwa takie
symbole odpowiada mnożeniu testu przez ich iloczyn.
Dla $s=2$ mamy $b_2=1+|\xi|^2$, toteż
$J^2=I-\Delta$ dzięki twierdzeniu
~\ref{thm:fourier-auto-tempered22}.

Dla $m\in\N_0$ zapis $W^{m,2}(\R^d)$ oznacza
przestrzeń funkcji $u\in L^2(\R^d)$, których
wszystkie pochodne słabe $\partial^\alpha u$,
$|\alpha|\le m$, należą do $L^2$.
Przyjmujemy normę
\[
 \|u\|_{W^{m,2}}^2=\sum_{|\alpha|\le m}
 \|\partial^\alpha u\|_2^2,
\]
włączając $\partial^0u=u$.

\begin{uwaga}[Szeregi Fouriera i transformata na $\R^d$]
W rozdziale \ref{sec:dt-series-fourier} na okręgu używamy
znormalizowanej miary Haara. Tutaj stosujemy miarę
Lebesgue'a i unitarną transformatę Fouriera z definicji
\ref{def:fourier-schwartz17}; w normie $H^s$ nie występuje
dodatkowy czynnik normalizujący.
\end{uwaga}

\begin{twierdzenie}[Podstawowe własności $H^s$]
\label{thm:dt-Hs}
Dla dowolnych $s,t\in\R$ zachodzą:
\begin{enumerate}
\item $H^s(\R^d)$ jest przestrzenią Hilberta, a
 $J^s:H^s\to L^2$ jest izometrycznym izomorfizmem
 na całą przestrzeń $L^2$.
\item Jeśli $s\ge t$, to $H^s\hookrightarrow H^t$
 ciągle i $\|u\|_{H^t}\le\|u\|_{H^s}$.
\item Dla $\alpha\in\N_0^d$, $\partial^\alpha:H^s\to
 H^{s-|\alpha|}$ jest ograniczony, z normą
 operatorową co najwyżej $1$.
\item Dla $m\in\N_0$ zachodzi $H^m(\R^d)=W^{m,2}(\R^d)$
 ze wzajemnie równoważnymi normami, po utożsamieniu
 funkcji z dystrybucjami regularnymi.
\item $\mathcal S(\R^d)$ jest gęsta w $H^s(\R^d)$.
 Ponadto $J^t:H^s\to H^{s-t}$ jest izometrycznym
 izomorfizmem na całą przestrzeń $H^{s-t}$.
\end{enumerate}
\end{twierdzenie}
\begin{proof}
Niech $u\in H^s$ i $b_s\widehat u=T_h$,
gdzie $h\in L^2$. Niech $v=\mathcal F^{-1}h$
w sensie $L^2$. Zgodność transformat na $L^2$
i $\mathcal S'$ z twierdzenia~\ref{thm:plancherel22}
daje $\mathcal FT_v=T_h$, a odwracalność
na $\mathcal S'$ daje $J^su=T_v$.
Wobec tego $J^s$ przyjmuje wartości w $L^2$
i z twierdzenia~\ref{thm:plancherel22}
otrzymujemy
\[
 \|J^s u\|_2=\|b_s\widehat u\|_2
 =\|u\|_{H^s}.
\]
Zapis normy pośrodku oznacza normę reprezentanta
$h$ z definicji. Dla dowolnego $v\in L^2$
traktujemy $v$ jako dystrybucję temperowaną
i określamy $u=J^{-s}v$. Wówczas
$b_s\widehat u=\widehat v\in L^2$, a $J^su=v$.
Iniektywność $J^s$ wynika z istnienia $J^{-s}$.
Iloczyn skalarny definiujemy przez
\[
 \langle u,v\rangle_{H^s}
 =\langle J^s u,J^s v\rangle_2.
\]
Liniowość $J^s$ i własności iloczynu skalarnego
na $L^2$ dają liniowość względem pierwszego
argumentu, antyliniowość względem drugiego
i symetrię hermitowską.
Ponadto $\langle u,u\rangle_{H^s}=\|J^su\|_2^2$
jest nieujemne, a równość z zerem daje $J^su=0$
i stąd $u=0$ dzięki iniektywności.
Iloczyn skalarny wyznacza podaną normę.
Jeśli $(u_n)$ jest ciągiem Cauchy'ego w $H^s$,
to $(J^su_n)$ jest ciągiem Cauchy'ego w $L^2$
i ma tam granicę $v$. Dla $u=J^{-s}v$ mamy
$\|u_n-u\|_{H^s}=\|J^s u_n-v\|_2\to0$.
Stąd zupełność i pierwsza teza.

Jeśli $s\ge t$, to $0<b_{t-s}\le1$ i
$b_t\widehat u=b_{t-s}(b_s\widehat u)$.
Prawa strona jest reprezentowana przez funkcję
$L^2$, a jej norma jest nie większa niż norma
$b_s\widehat u$. Daje to drugą tezę.

Dla dowolnego wielowskaźnika $\alpha$ mamy
na mocy twierdzenia~\ref{thm:fourier-auto-tempered22}
\[
 b_{s-|\alpha|}\mathcal F(\partial^\alpha u)
 =b_{-|\alpha|}(i\xi)^\alpha b_s\widehat u.
\]
Moduł mnożnika po prawej stronie nie przekracza
$1$, ponieważ
$|\xi^\alpha|\le(1+|\xi|^2)^{|\alpha|/2}$.
Zatem $\partial^\alpha u\in H^{s-|\alpha|}$
i $\|\partial^\alpha u\|_{H^{s-|\alpha|}}
\le\|u\|_{H^s}$.

Niech teraz $m\in\N_0$. Z rozwinięcia
wielomianowego otrzymujemy
\[
 (1+|\xi|^2)^m
 =\sum_{|\alpha|\le m}c_\alpha|\xi^\alpha|^2,
 \qquad c_\alpha=\frac{m!}{(m-|\alpha|)!\alpha!}.
\]
Każde $c_\alpha$ jest dodatnią liczbą całkowitą,
a z podstawienia $\xi=(1,\ldots,1)$ wynika
$\sum_{|\alpha|\le m}c_\alpha=(d+1)^m$.
W szczególności
\[
 \sum_{|\alpha|\le m}|\xi^\alpha|^2
 \le (1+|\xi|^2)^m
 \le (d+1)^m\sum_{|\alpha|\le m}|\xi^\alpha|^2.
\]
Jeśli $u\in H^m$, to z drugiej i trzeciej tezy
wynika, że $\partial^\alpha u\in H^{m-|\alpha|}$
dla $|\alpha|\le m$ (włącznie z $\alpha=0$).
Ponieważ $m-|\alpha|\ge0$, drugą tezę
stosujemy z wykładnikami $m-|\alpha|$ i $0$.
Zatem wszystkie te pochodne należą do $H^0=L^2$.
Utożsamienie $H^0=L^2$ wynika z pierwszej tezy
dla $s=0$, ponieważ $J^0=I$.
Po ograniczeniu do testów o zwartym nośniku
pochodne dystrybucyjne reprezentowane przez
funkcje $L^2$ są pochodnymi słabymi w sensie
rozdziału~\ref{sec:pochodne20}.
A zatem $u\in W^{m,2}$.

Odwrotnie, dla $u\in W^{m,2}$ każda słaba
pochodna $v_\alpha=\partial^\alpha u$ należy
do $L^2$. Równość dystrybucyjna
$\partial^\alpha T_u=T_{v_\alpha}$ zachodzi
na $\mathcal D$, więc zachodzi też na $\mathcal S$:
obie strony są temperowane, a ograniczenie
$\mathcal S'\to\mathcal D'$ jest iniektywne.
Po transformacji dostajemy
$(i\xi)^\alpha\widehat u=\widehat v_\alpha$
jako dystrybucje temperowane.
Lewą stronę reprezentuje funkcja lokalnie
całkowalna, ponieważ $\widehat u\in L^2$
i wielomian jest ograniczony na zwartych zbiorach.
Z iniektywności dystrybucji regularnych
równość obowiązuje prawie wszędzie.
Stąd z Plancherela i skończonej sumy
\[
 \int_{\R^d}
 \left(\sum_{|\alpha|\le m}|\xi^\alpha|^2\right)
 |\widehat u(\xi)|^2\,\mathrm d\lambda_d(\xi)
 =\sum_{|\alpha|\le m}\|v_\alpha\|_2^2<\infty.
\]
Porównanie wielomianów daje $u\in H^m$.
Ten sam rachunek obowiązuje dla $u\in H^m$,
ponieważ wykazaliśmy już przynależność
do $W^{m,2}$. W obu kierunkach otrzymujemy
\[
 \|u\|_{W^{m,2}}\le\|u\|_{H^m}
 \le(d+1)^{m/2}\|u\|_{W^{m,2}}.
\]

Na koniec $J^s$ i $J^{-s}$ odwzorowują
$\mathcal S$ w siebie, ponieważ transformata
i odwrotność zachowują $\mathcal S$,
a symbole $b_s,b_{-s}$ są dopuszczalnymi
mnożnikami. Każda funkcja Schwartza należy
do $H^s$, gdyż $b_s\widehat u\in\mathcal S\subset L^2$.
Dla $u\in H^s$ wybieramy $v_n\in\mathcal S$
zbiegające do $J^su$ w $L^2$ na mocy
gęstości $\mathcal S$ w $L^2$ wykorzystanej w dowodzie
twierdzenia~\ref{thm:plancherel22}. Wtedy
$u_n=J^{-s}v_n\in\mathcal S$ i
$\|u_n-u\|_{H^s}=\|v_n-J^su\|_2\to0$.
Dla dowolnego $t$ mamy ponadto
$b_{s-t}\mathcal F(J^tu)=b_s\widehat u$,
więc $\|J^tu\|_{H^{s-t}}=\|u\|_{H^s}$.
Dla $v\in H^{s-t}$ stosujemy tę samą tożsamość
z wykładnikami $s-t$ i $-t$, otrzymując
$J^{-t}v\in H^s$.
Ponieważ $J^tJ^{-t}v=v$ i $J^{-t}J^tu=u$,
operator $J^{-t}$ jest odwrotnością na całym
$H^{s-t}$.
\end{proof}

\begin{wniosek}[Kiedy delta należy do ujemnej przestrzeni Sobolewa?]
\label{cor:dt-delta-Hs}
Dla $s\in\R$ zachodzi
\[
 \delta_0\in H^s(\R^d)\quad\Longleftrightarrow\quad
 s<-\frac d2.
\]
\end{wniosek}
\begin{proof}
Z przykładu~\ref{ex:fourier-delta22} mamy
$\widehat{\delta_0}=(2\pi)^{-d/2}$ jako dystrybucję regularną.
Wobec tego $b_s\widehat{\delta_0}$ jest
reprezentowana przez funkcję $(2\pi)^{-d/2}b_s$.
Stały niezerowy czynnik nie wpływa na warunek
kwadratowalności.
Z iniektywności dystrybucji regularnych
warunek przynależności do $H^s$ jest równoważny
\[
 \int_{\R^d}(1+|\xi|^2)^s\,
 \mathrm d\lambda_d(\xi)<\infty.
\]
Na domkniętej kuli jednostkowej funkcja podcałkowa
jest ciągła i ograniczona. Dla $r\ge1$ zachodzi
\[
 \min\{1,2^s\}r^{2s}\le(1+r^2)^s
 \le\max\{1,2^s\}r^{2s}.
\]
Dla $d\ge2$ całkowanie radialne poza kulą
sprowadza zatem warunek do skończoności
$\int_1^\infty r^{2s+d-1}\,\mathrm d\lambda_1(r)$;
dla $d=1$ całkujemy osobno na obu półprostych
i otrzymujemy ten sam warunek z $d=1$.
Całka potęgowa jest skończona dokładnie wtedy,
gdy $2s+d-1<-1$, czyli $s<-d/2$.
Dla $s=-d/2$ jest to całka z $r^{-1}$,
rozbieżna logarytmicznie. W szczególności
$H^s$ dla $s<-d/2$ zawiera dystrybucję
niereprezentowaną przez funkcję lokalnie całkowalną:
nieregularność $\delta_0$ wykazano w przykładzie
~\ref{ex:delta19}.
\end{proof}


\subsection{Zasada nieoznaczoności}
\label{sec:dt-uncertainty}
\providecommand{\atlasref}[2]{#2}

Por. \atlasref{PDE-08::thm:pde8-uncertainty}{RRC-08: nierówność Heisenberga}.


\begin{twierdzenie}[Nierówność Heisenberga w jednej współrzędnej]
\label{thm:dt-uncertainty}
Jeśli $0\ne f\in\mathcal S(\R^d)$ i $1\le j\le d$, to
\[
 \|x_jf\|_{L^2}\,\|\xi_j\widehat f\|_{L^2}
 \ge\frac12\,\|f\|_{L^2}\,\|\widehat f\|_{L^2}.
\]
Stała $1/2$ jest optymalna. Równość zachodzi
dokładnie wtedy, gdy dla pewnego $\lambda>0$
funkcja $f$ ma postać
\[
 f(x)=e^{-\lambda x_j^2/2}\,g(x'),
\]
gdzie dla $d>1$ przez $x'$ oznaczamy pozostałe
$d-1$ współrzędnych, a $g\in\mathcal S(\R^{d-1})$
jest niezerowa. Dla $d=1$ czynnik $g(x')$
zastępujemy niezerową stałą zespoloną.
\end{twierdzenie}
\begin{proof}
Najpierw zauważmy, że $|f|^2=f\overline f$,
a także $x_j|f|^2$ należą do $\mathcal S$.
Sprzężenie zespolone nie zmienia półnorm.
Dla $u,v\in\mathcal S$ wzór Leibniza daje
\[
 Q_{N,m}(uv)\le 2^m Q_{N,m}(u)Q_{0,m}(v)
 \qquad(N,m\in\N_0),
\]
gdyż każda pochodna iloczynu jest sumą składników
$(\partial^\gamma u)(\partial^{\alpha-\gamma}v)$
ze współczynnikami dwumianowymi o sumie $2^{|\alpha|}\le2^m$.
Iloczyn dwóch funkcji Schwartza należy zatem do
$\mathcal S$, a mnożenie przez współrzędną zachowuje
$\mathcal S$ na mocy stwierdzenia~\ref{prop:schwartz-operators17}.
W szczególności $\partial_j(x_j|f|^2)\in L^1$.
Dla ustalonych pozostałych współrzędnych
podstawowe twierdzenie rachunku całkowego daje
\[
 \int_{-R}^R\partial_j(x_j|f(x)|^2)
 \,\mathrm d\lambda_1(x_j)
 =\left[x_j|f(x)|^2\right]_{x_j=-R}^{x_j=R}.
\]
Prawa strona dąży do zera dla $R\to\infty$,
ponieważ dla $|x_j|=R$ mamy
\[
 |x_j||f(x)|^2
 \le R\,q_{1,0}(f)^2(1+R)^{-2}.
\]
Dla $d>1$ dzięki całkowalności pochodnej
na $\R^d$ twierdzenie Fubiniego zapewnia, że
dla prawie każdego $x'$ pochodna ta jest
bezwzględnie całkowalna na prostej zmiennej $x_j$.
Na takim przekroju całka po $[-R,R]$ zbiega
do całki po $\R$ z twierdzenia o zbieżności
zmajoryzowanej, z majorantą
$|\partial_j(x_j|f|^2)|$ na przekroju.
Z oszacowania składnika brzegowego wynika,
że ta całka po $\R$ jest zerowa. Fubini
pozwala teraz całkować tę równość względem $x'$.
Dla $d=1$ przechodzimy do granicy bezpośrednio,
ponieważ pochodna jest w $L^1(\R)$.
Otrzymujemy
\[
 \int_{\R^d}\partial_j(x_j|f|^2)
 \,\mathrm d\lambda_d(x)=0.
\]
Z reguły Leibniza
\[
 \partial_j(x_j|f|^2)
 =|f|^2+2\operatorname{Re}(x_j\overline f\,\partial_j f).
\]
Oba składniki są całkowalne: dla drugiego
wynika to z nierówności Cauchy'ego--Schwarza
dla $x_jf,\partial_j f\in L^2$.
W konsekwencji
\[
 \begin{aligned}
 \|f\|_2^2
 &=-2\operatorname{Re}
 \int_{\R^d}x_j\overline f(x)\,\partial_j f(x)
 \,\mathrm d\lambda_d(x)\\
 &\le 2\left|
 \int_{\R^d}x_j\overline f(x)\,\partial_j f(x)
 \,\mathrm d\lambda_d(x)\right|\\
 &\le 2\|x_j f\|_2\|\partial_j f\|_2.
 \end{aligned}
\]
Z twierdzeń~\ref{thm:fourier-schwartz17}
i \ref{thm:plancherel22} mamy
\[
 \|\partial_j f\|_2=\|\xi_j\widehat f\|_2,
 \qquad
 \|f\|_2=\|\widehat f\|_2.
\]
Po podstawieniu otrzymujemy
$\|f\|_2^2\le2\|x_jf\|_2\|\xi_j\widehat f\|_2$.
Dzieląc przez $2$ i wykorzystując zachowanie normy
przez unitarną transformatę, otrzymujemy tezę.

Ustalmy teraz warunek równości. Ponieważ
$\|f\|_2>0$, powyższe oszacowanie pokazuje,
że zarówno $x_jf$, jak i $\partial_j f$
mają dodatnie normy $L^2$.
Równość w nierówności Heisenberga jest
równoważna równości w obu nierównościach
w powyższym łańcuchu.
Równość w nierówności Cauchy'ego--Schwarza
oznacza liniową zależność obu wektorów w $L^2$:
$\partial_j f=cx_jf$ dla pewnego $c\in\C$.
Istotnie, dla $p=\partial_j f$, $q=x_jf$
oraz $c=\langle p,q\rangle_2/\|q\|_2^2$
mamy
\[
 \|p-cq\|_2^2
 =\|p\|_2^2-\frac{|\langle p,q\rangle_2|^2}{\|q\|_2^2}.
\]
Równość w nierówności Cauchy'ego--Schwarza
daje zatem $p=cq$ w $L^2$, a zależność ta
jest również wystarczająca dla równości.
Wtedy
\[
 \int_{\R^d}x_j\overline f\,\partial_j f
 \,\mathrm d\lambda_d(x)=c\|x_jf\|_2^2.
\]
Równość w pierwszej nierówności wymaga,
by liczba po lewej stronie była rzeczywista
i ujemna (jej część rzeczywista wynosi
$-\|f\|_2^2/2$). Stąd $c=-\lambda$
dla pewnego $\lambda>0$.
Ponieważ $\partial_j f$ i $x_j f$ są ciągłe,
ich równość prawie wszędzie daje równość
w każdym punkcie: niezerowa różnica w punkcie
byłaby niezerowa na pewnej otwartej kuli
o dodatniej mierze.
Dla ustalonych $x'$ otrzymujemy więc
\[
 \partial_j\bigl(e^{\lambda x_j^2/2}f(x)\bigr)=0,
\]
czyli $f(x)=e^{-\lambda x_j^2/2}g(x')$,
gdzie $g(x')=f(x)|_{x_j=0}$.
Dla $d>1$ funkcja $g$ jest gładka jako
ograniczenie funkcji gładkiej do hiperpłaszczyzny.
Jeśli $\alpha',\beta'\in\N_0^{d-1}$,
to przez $\alpha,\beta\in\N_0^d$ oznaczmy
wielowskaźniki powstałe przez wstawienie
zera na $j$-tej pozycji.
Wtedy
\[
 \sup_{x'\in\R^{d-1}}
 |(x')^{\alpha'}\partial^{\beta'}g(x')|
 \le p_{\alpha,\beta}(f)<\infty.
\]
Zatem $g\in\mathcal S(\R^{d-1})$ i $g\ne0$,
bo w przeciwnym razie wzór na $f$ dawałby
$f=0$. Dla $d=1$ tym samym argumentem
otrzymujemy niezerową stałą $f(0)$.

Odwrotnie, dla dowolnej $\lambda>0$ i takiej
niezerowej $g$ funkcja
$f(x)=e^{-\lambda x_j^2/2}g(x')$ należy
do $\mathcal S$. Dla $d>1$, oznaczając
$G(t)=e^{-\lambda t^2/2}$, mamy
\[
 p_{\alpha,\beta}(f)=
 p_{\alpha_j,\beta_j}(G)\,
 p_{\alpha',\beta'}(g)<\infty,
\]
gdzie $\alpha',\beta'$ powstają przez
usunięcie $j$-tej współrzędnej.
Wynika to z rozdzielenia pochodnych i wag
oraz supremum niezależnych nieujemnych czynników.
Dla $d=1$ wystarcza przynależność Gaussa do $\mathcal S$.
Niezerowość $f$ wynika z niezerowości $g$.
Spełnia $\partial_j f=-\lambda x_j f$,
więc obie nierówności w dowodzie stają się
równościami. Takie funkcje istnieją w każdym
wymiarze, na przykład $f(x)=e^{-|x|^2/2}$.
Zatem stałej $1/2$ nie można zwiększyć.
\end{proof}

\begin{wniosek}[Nierówność z przesuniętymi środkami]
Dla $f\in\mathcal S(\R^d)$ i dowolnych $a,b\in\R$
mamy
\[
 \|(x_j-a)f\|_2\,\|(\xi_j-b)\widehat f\|_2
 \ge\tfrac12\|f\|_2\|\widehat f\|_2.
\]
\end{wniosek}
\begin{proof}
Dla $e_j$ oznaczającego $j$-ty wektor
bazy standardowej połóżmy
$h(x)=e^{-ibx_j}f(x+ae_j)$.
Przesunięcie zachowuje $\mathcal S$ przez
nierówność $1+|x|\le(1+|a|)(1+|x+ae_j|)$,
a mnożenie przez $e^{-ibx_j}$ zachowuje
$\mathcal S$, gdyż wszystkie pochodne
mnożnika są ograniczone.
Ponieważ $h\in\mathcal S\subset L^1$,
podstawienie $y=x+ae_j$ w całce Fouriera daje
\[
 \begin{aligned}
 \widehat h(\xi)
 &=(2\pi)^{-d/2}\int_{\R^d}e^{-ix\cdot(\xi+be_j)}
 f(x+ae_j)\,\mathrm d\lambda_d(x)\\
 &=e^{ia(\xi_j+b)}(2\pi)^{-d/2}
 \int_{\R^d}e^{-iy\cdot(\xi+be_j)}f(y)
 \,\mathrm d\lambda_d(y)\\
 &=e^{ia(\xi_j+b)}\widehat f(\xi+be_j).
 \end{aligned}
\]
Czynniki wykładnicze mają moduł $1$,
a miara Lebesgue'a jest niezmiennicza
względem przesunięć. Stąd
\[
 \begin{aligned}
 \|h\|_2&=\|f\|_2,&
 \|x_jh\|_2&=\|(x_j-a)f\|_2,\\
 \|\widehat h\|_2&=\|\widehat f\|_2,&
 \|\xi_j\widehat h\|_2&=\|(\xi_j-b)\widehat f\|_2.
 \end{aligned}
\]
Dla $f\ne0$ stosujemy twierdzenie
~\ref{thm:dt-uncertainty} do $h$;
dla $f=0$ nierówność jest oczywista.
\end{proof}

Dla $f\ne0$ można w szczególności wybrać
\[
 a=\frac{\int_{\R^d}x_j|f(x)|^2\,
 \mathrm d\lambda_d(x)}{\|f\|_2^2},
 \qquad
 b=\frac{\int_{\R^d}\xi_j|\widehat f(\xi)|^2\,
 \mathrm d\lambda_d(\xi)}{\|\widehat f\|_2^2}.
\]
Całki są bezwzględnie zbieżne, ponieważ
\[
 \begin{aligned}
 \int_{\R^d}|x_j||f(x)|^2\,\mathrm d\lambda_d(x)
 &\le\|x_jf\|_2\|f\|_2,\\
 \int_{\R^d}|\xi_j||\widehat f(\xi)|^2\,
 \mathrm d\lambda_d(\xi)
 &\le\|\xi_j\widehat f\|_2\|\widehat f\|_2.
 \end{aligned}
\]
Oba mianowniki są dodatnie z niezerowości $f$
i twierdzenia Plancherela. Liczby $a,b$ są
rzeczywiste, a obie niżej wymienione gęstości
mają całkę $1$ z definicji normy $L^2$.
Drugie momenty są skończone, ponieważ
$x_jf,\xi_j\widehat f\in L^2$.
Jeśli $\sigma_x^2$ i $\sigma_\xi^2$
oznaczają wariancje $j$-tej współrzędnej
względem miar probabilistycznych o gęstościach
$|f|^2/\|f\|_2^2$ oraz
$|\widehat f|^2/\|\widehat f\|_2^2$, to
\[
 \sigma_x=\frac{\|(x_j-a)f\|_2}{\|f\|_2},
 \qquad
 \sigma_\xi=\frac{\|(\xi_j-b)\widehat f\|_2}
 {\|\widehat f\|_2}.
\]
Z poprzedniego wniosku dostajemy
$\sigma_x\sigma_\xi\ge1/2$.
Stała odpowiada konwencji z definicji~\ref{def:fourier-schwartz17},
w której wykładnik transformaty to $-ix\cdot\xi$.
